\chapter*{Agradecimentos}

% \vspace{15mm}

Em primeiro lugar, dedico todo este trabalho ao Nosso Senhor Jesus Cristo, Rei dos céus e da terra. Nâo menos também deposito minhas gratidões ao sobrenatural título de Nossa Senhora de Fátima, razão pela qual me despertou o interesse pela busca constante da verdade.

Ademais, o fato consumado que é ser criado bacharel não seria possível sem incondicional apoio de minha família, a quem secundariamente devo o dom da vida. Não obstante, devo minhas considerações a Rogério, Vânia e seu filho João Pedro, do qual tenho uma eterna amizade, por cederem o lar que deu o estopim desse sonho.

Ao longo de toda a graduação os momentos de fortaleza foram estabelecidos pelos Nerdolas, a quem devo todos os bons momentos. Por isso, minha gratidão à Diogo Tallys, Filipe Soriano, Hugo Coelho, Isaque Cerqueira, José Matheus, Lucas Heron, Rafael Luciano, Ruan Tenório, Thiago Ribeiro, Victor Hugo e Vinícius Neitzke. 

Destarte, o rumo racional e espiritual que me fez presente durante esta jornada não seria possível sem as devidas providências da Legião de Maria, a quem eternamente serei grato pelas companhias de Danilo, Euricledes, João Pedro, Juliana e Klyvia. Outrossim, destaco aqui a fundamental importância dos coelgas que me guiaram no caminho da razão: Ismael Damião, Giselle Medeiros, Letícia Feitosa e Rômulo Souza.

Ainda devo minhas gratificações ao GEMA, do qual obtive muitas boas amizades com João Ayalla, Davi Romão, Laura Mainero, Davi Celestino, Gabriel Seixas, Beatriz Rodrigues, Ricardo, e Davi Vieira.


\begin{minipage}{.45\linewidth}
\begin{flushleft}
% Maceió, 24 de novembro de 2023.
\end{flushleft}
\end{minipage}
\begin{minipage}{.45\linewidth}
\end{minipage}